% ==============================================================================
% tabletop.tex — UE SIS589
% Guide de l'Animateur — Exercice Tabletop Ransomware LiZbiZ-SIS
% ==============================================================================

\chapter*{Guide de l'Animateur~— Tabletop Ransomware}
\addcontentsline{toc}{chapter}{Guide de l'Animateur: Tabletop}
\label{ann:tabletop}

\begin{boxinfo}
\textbf{Document réservé à l'enseignant.} Ce guide est distribué uniquement
à l'animateur du Tabletop. Il ne doit pas être communiqué aux équipes avant
la clôture de l'exercice.
\end{boxinfo}

% ==============================================================================
\section*{Objectifs pédagogiques de l'exercice}
% ==============================================================================

Cet exercice Tabletop teste la capacité des équipes à~:
\begin{itemize}
  \item Prendre des décisions sous pression et incertitude.
  \item Appliquer le cycle ISO~27035 sans assistance documentaire.
  \item Articuler les dimensions technique, légale et managériale
        d'une crise cyber.
  \item Communiquer de manière adaptée à différentes audiences.
  \item Identifier spontanément les lacunes de préparation.
\end{itemize}

% ==============================================================================
\section*{Chronologie des injects}
% ==============================================================================

\begin{table}[H]
\centering\small
\caption{Planning des injects — exercice de 4~heures}
\label{tab:tabletop-planning}
\rowcolors{2}{TableauPair}{TableauImpair}
\begin{tabular}{p{2cm}p{3cm}p{8.5cm}}
\toprule
\tableauHeader{Heure} & \tableauHeader{Événement} &
\tableauHeader{Objectif pédagogique} \\
\midrule
T+0   & Scénario initial & Triage, qualification, premières décisions \\
T+15  & Inject~1 (RDP + exfil) & Révision de l'évaluation, vecteur initial \\
T+30  & Bilan Phase~1 & Vérifier que les équipes ont activé l'astreinte \\
T+35  & Inject~2 (ransomware confirmé) & Confinement, mapping MITRE, comm DG \\
T+75  & Inject~3 (backups chiffrés + MoMo) & Escalade légale, décision restauration \\
T+90  & Bilan Phase~2 & Vérifier cohérence des décisions \\
T+120 & Phase~3 (éradication) & Ordonnancement des actions \\
T+150 & Début RETEX & Réflexivité, auto-évaluation \\
T+180 & Présentations orales & Communication, maîtrise du sujet \\
T+210 & Débriefing animateur & Retour sur les décisions clés \\
\bottomrule
\end{tabular}
\end{table}

% ==============================================================================
\section*{Éléments de réponse attendus}
% ==============================================================================

\subsection*{Phase~1 — Décisions attendues}

\textbf{Classification correcte~:} Gravité \gravite{critique}
(données clients, service de paiement, propagation probable).
Activation immédiate~: RSSI (même en déplacement), Incident Manager,
analyste N2.

\textbf{5~premières actions (ordre attendu)~:}
\begin{enumerate}
  \item Capturer la mémoire de srv-erp-01 AVANT toute action (avml ou WinPmem).
  \item Isoler srv-erp-01 du réseau (VLAN quarantaine, pas d'extinction).
  \item Bloquer l'IP sortante 185.220.101.5 au firewall.
  \item Désactiver le compte \texttt{administrator} sur srv-erp-01.
  \item Notifier le RSSI et le DG.
\end{enumerate}

\textbf{Points de vigilance~:} Équipes qui éteignent srv-erp-01 immédiatement
sans capturer la RAM~: pénalité pédagogique~; rappeler l'ordre de volatilité.
Équipes qui ne pensent pas à isoler srv-backup-01~: rappeler la propagation SMB.

\subsection*{Inject~1 — Réponse attendue}

Le vecteur initial est l'accès RDP depuis 196.207.12.88 le 14~mai à 22h41
(veille de l'incident). L'exfiltration vers 185.220.101.5 s'est faite
pendant 7h34 avant le début du chiffrement.
\textbf{TTPs attendus~:} T1133 (External Remote Services),
T1078 (Valid Accounts), T1071.001 (Application Layer Protocol),
T1048 (Exfiltration over Alternative Protocol).

\subsection*{Phase~2 — Réponse à la recommandation de paiement}

\textbf{Recommandation académique~: Ne pas payer.}
Arguments~:
\begin{enumerate}
  \item Aucune garantie de déchiffrement~; 30\% des victimes qui paient
        ne récupèrent pas leurs données.
  \item Finance la criminalité et attire d'autres attaques.
  \item Potentiellement illégal si le groupe est sanctionné (OFAC).
  \item La clé de déchiffrement peut elle-même contenir un backdoor.
  \item Une sauvegarde tape hors-site existe (inject~3).
\end{enumerate}

\subsection*{Inject~3 — Obligations légales}

Les transactions MTN MoMo suspectes impliquent une obligation de
notification à MTN, à l'ANTIC et potentiellement au BEAC
(Banque des États de l'Afrique Centrale) selon les montants.
La fuite de données clients déclenche les obligations de la Loi~2010/012.

% ==============================================================================
\section*{Grille d'observation de l'animateur}
% ==============================================================================

\begin{table}[H]
\centering\small
\caption{Grille d'observation — comportements à noter}
\label{tab:tabletop-observation}
\rowcolors{2}{TableauPair}{TableauImpair}
\begin{tabular}{p{6cm}p{3cm}p{3cm}p{1.8cm}}
\toprule
\tableauHeader{Comportement observé} & \tableauHeader{Moment} &
\tableauHeader{Équipe} & \tableauHeader{Score} \\
\midrule
A capturé la RAM avant isolation & Phase~1 & & /2 \\
A isolé correctement srv-erp-01 & Phase~1 & & /2 \\
A pensé à srv-backup-01 spontanément & Phase~1 & & /1 \\
Communication interne adaptée & Phase~1 & & /1 \\
Mapping MITRE complet ($\geq 8$~TTPs) & Phase~2 & & /2 \\
Recommandation paiement justifiée & Phase~2 & & /2 \\
A identifié les obligations légales MoMo & Phase~2--3 & & /2 \\
Plan d'éradication ordonné correctement & Phase~3 & & /2 \\
RETEX sincère (lacunes reconnues) & RETEX & & /2 \\
Présentation orale structurée & Oral & & /2 \\
\bottomrule
\end{tabular}
\end{table}

% ==============================================================================
\section*{Questions de débriefing (animateur)}
% ==============================================================================

\begin{enumerate}
  \item \og{}Pourquoi n'avez-vous pas capturé la mémoire vive en premier~?\fg{}
        (si non fait)
  \item \og{}Quelle était votre source d'information principale pour prendre
        la décision de confinement~? Était-elle suffisante~?\fg{}
  \item \og{}Si vous aviez eu un playbook ransomware, qu'aurait-il contenu
        que vous avez découvert en cours de simulation~?\fg{}
  \item \og{}La communication vers le DG était-elle adaptée~?
        Trop technique~? Trop floue~?\fg{}
  \item \og{}Avez-vous pensé à la dimension légale (ANTIC, obligations
        de notification) dès le départ ou seulement après l'inject~3~?\fg{}
  \item \og{}Que feriez-vous différemment si cet incident se reproduisait
        dans 6~mois~? Quels artefacts documentaires manquaient~?\fg{}
\end{enumerate}
